\documentclass[12pt,a4paper]{article}
\usepackage{geometry}
\usepackage{amsmath,amssymb}
\usepackage{graphicx}
\usepackage{booktabs}
\usepackage{hyperref}
\usepackage{cite}
\usepackage{setspace}
\usepackage{caption}
\usepackage{array}
\geometry{margin=1in}
\onehalfspacing

\hypersetup{
  colorlinks=true,
  linkcolor=blue,
  citecolor=blue,
  urlcolor=blue
}

\title{
  \textbf{BIMOD-006} \\[0.5em]
  \Large BMNP Meridian Hypothesis: \\
  Cross-Polar Single-Domain Nanoparticle Network, \\
  Optical Gating, and Eight-Constitutional Classification --- \\
  An Integrated Theory
}
\author{
  DuckJin Yoon (Solvaram) \\
  \textit{Eight-Constitutional Research Platform} \\
  \small with collaborative input from Claude AI (Anthropic) \\
  \small \texttt{OSF: https://doi.org/10.17605/OSF.IO/VJF63}
}
\date{Version 1.3 --- February 2026}

\begin{document}
\maketitle

% ─────────────────────────────────────────
\begin{abstract}
BIMOD-006 proposes the \textbf{BMNP Meridian Hypothesis}: biogenic magnetic
nanoparticles (BMNPs) in living organisms are each in a stable
\emph{single magnetic domain} (SMD) state, exhibiting unambiguous
single-polarity behavior.  These SMD particles organize into two
complementary structures: (1)~\emph{cross-polar alternating chains}
(N--S--N--S\ldots) serving as directional signal conduits, and
(2)~\emph{monopolar clusters} at nodal acupoints acting as functional
terminals.  Together they constitute the proposed physical substrate of
the classical meridian system of East Asian medicine.

Key lines of evidence are: (i)~darkened magnetosome particles in
\textit{Gammaproteobacteria} SS-5 TEM images consistently occupy
S-polarity positions (novel BIMOD observation); (ii)~each magnetosome is
confirmed as an independent SMD nanoparticle
(\textit{Nanoscale Advances}, 2020); (iii)~BMNP chains are detected in
plant phloem sieve-tube cell walls; (iv)~injected magnetic nanoparticles
travel along classical meridian pathways \cite{johng2007}; (v)~BIMOD
polarity discrimination is photon-gated.

Eight constitutional types expressed as four bilateral Nasion polarity
phenotypes are proposed to reflect eight genetically determined BMNP
network architectures, each with a unique organ dominance sequence.
The framework yields falsifiable predictions and provides a foundation
for developing objective sensor systems.

\textbf{Keywords:} BIMOD, biogenic magnetic nanoparticles, meridian
hypothesis, cross-polar chain, single magnetic domain, monopolar node,
eight constitution, Nasion acupoint, optical gating, magnetosome,
constitutional pharmacology
\end{abstract}

% ─────────────────────────────────────────
\tableofcontents
\newpage

% ─────────────────────────────────────────
\section{Introduction}

\subsection{The Unresolved Problem of Meridians}
The meridian system of classical East Asian medicine has resisted
physical characterization despite millennia of clinical application.
Proposed substrates---connective-tissue planes, interstitial fluid
channels, neural pathways, and the Bonghan duct system---have each
failed to achieve consensus.

We propose a new candidate: a network of BMNPs in which each particle
occupies a stable SMD state with unambiguous single polarity, arranged
in N--S--N--S alternating chains along anatomically stable pathways,
with monopolar clusters at nodal points.

\subsection{Single-Domain Nanoparticles and the Monopolar Phenomenon}
A magnetic nanoparticle exceeding the critical size for stable SMD
formation ($\sim$25~nm) exhibits a single-polarity character: all
internal spins align in one direction.  Published magnetosome research
\cite{kornig2020} confirms that each particle is an independent SMD
nanoparticle with uniaxial magnetic anisotropy along the chain axis;
magnetization reversal occurs only by coherent rotation, not by physical
displacement.

BIMOD observations are consistent with this: in permanent-magnet
cross-sections the central $\sim$30\% zone shows the primary polarity,
the outermost rim shows a thin band of opposite polarity, and the
intermediate zone is non-responsive---a classic dipole signature.
By contrast, nanoparticle specimens yield a \emph{uniform single-polarity
response across the entire area}, with no intermediate non-polar zone
and no opposite-polarity rim.  This constitutes direct BIMOD evidence
for SMD single-polarity behavior in BMNPs.

\subsection{The BIMOD Detection Framework}
BIMOD (Bodily Interferometric Magneto-Optical Discrimination) operates
through a quadruple resonance: (1)~a permanent-magnet probe, (2)~BMNP-
containing human tissue, (3)~Earth's geomagnetic field, and (4)~photon-
mediated optical gating.  Reproducible polarity discrimination has been
demonstrated from physical magnets, photographs, historical oil
portraits, and X-ray/MRI images, establishing structured optical
information as a polarity carrier.

\subsection{Key Observation: Darkened Particles Co-Localise with S-Polarity}
BIMOD analysis of published TEM images of magnetosomes in
\textit{Gammaproteobacteria} SS-5 (Wikipedia: \textit{Magnetoreception})
revealed strict N--S--N--S alternation across all three panels (A, B,
C).  Critically, every darkened/electron-dense particle fell at an
S-polarity position without exception---suggesting a coupling between
oxidation state and magnetic polarity identity.

% ─────────────────────────────────────────
\section{Theoretical Framework: The BMNP Meridian Hypothesis}

\subsection{Structural Components}

\begin{table}[h]
\centering
\caption{Proposed network components}
\begin{tabular}{>{\bfseries}p{3.2cm} p{5.5cm} p{5.5cm}}
\toprule
Component & Structure & Proposed function \\
\midrule
Cross-polar chain &
  N and S SMD particles in strict N--S--N--S sequence along
  anatomical pathway &
  Directional signal conduction via in-situ polarity relay;
  no particle displacement required \\[4pt]
Monopolar node (S) &
  S-polarity SMD cluster at acupoint locus &
  Signal termination, amplification, organ-specific transduction \\[4pt]
Monopolar node (N) &
  N-polarity SMD cluster at acupoint locus &
  Signal reception, constitutional baseline maintenance \\[4pt]
Nasion surface expression (4 bilateral loci, each N or S) &
  Four paired symmetric SMD nodes exposed at the cranial surface
  midline &
  Eight-constitutional phenotype identification via BIMOD \\
\bottomrule
\end{tabular}
\end{table}

\subsection{Signal Transmission Mechanism}
Polarity signals are transmitted without particle migration.
Two non-exclusive mechanisms are considered:

\paragraph{Possibility A --- Cross-polar serial relay.}
SMD particles fixed in tissue transmit polarity state changes along the
chain through nearest-neighbour magnetic interactions, analogous to
magnon (spin-wave) propagation.  Whether the interaction is strictly
monopolar or involves residual dipolar coupling requires further
investigation.

\paragraph{Possibility B --- Sequential redox cascade.}
The preferentially oxidised state of S-polarity nodes (Section~4)
constitutes an electron-transfer chain whose output is ion-channel
modulation.  This model supports in-situ signal transmission without
particle displacement.

The two models are not mutually exclusive: Possibility~A may govern
fast signal propagation while Possibility~B provides sustained
regulation.  Discrimination requires real-time SQUID-based measurements
along acupoint pathways.

\subsection{Eight-Constitutional Genotype Hypothesis}
The Nasion acupoint comprises \textbf{eight positions} distributed
symmetrically left and right of the facial midline.
Two positions of the same colour (see Figure~\ref{fig:nasionmap})
form one constitutional axis; left and right polarity are always
identical within the same constitution --- for example, on Axis~1,
both left and right positions are N for Pulmotonia and both S for
Hepatonia.  The facial midline itself is non-polar.
Because adjacent positions on either side may belong to \emph{different}
axes, the eight-position description avoids the misreading that two
neighbouring same-side positions share one axis.

The eight Nasion polarity phenotypes reflect eight genetically determined
BMNP network architectures.  Official constitutional nomenclature follows
the Eight-Constitution Medical Research Institute
(\url{www.ecmed.org}).\footnote{%
  Latin and Korean constitutional names, and organ dominance sequences,
  follow the official nomenclature of the Eight-Constitution Medical
  Research Institute (\url{www.ecmed.org}).}

\begin{table}[h]
\centering
\caption{Eight Nasion polarity phenotypes}
\begin{tabular}{llll}
\toprule
Nasion code & Latin name (official) & Korean name (official) &
  Dominant organ \\
\midrule
1S & Hepatonia   & MokyangcheJil & Liver (max) \\
1N & Pulmotonia  & GeumyangcheJil & Lung (max) \\
2N & Cholecystonia & MokeumcheJil & Gallbladder (max) \\
2S & Colonotonia & GeumeumcheJil & Colon (max) \\
3N & Gastrotonia & ToyangcheJil & Stomach (max) \\
3S & Vesicotonia & SuyangcheJil & Bladder (max) \\
4N & Pancreotonia & ToeumcheJil & Pancreas (max) \\
4S & Renotonia   & SueumcheJil & Kidney (max) \\
\bottomrule
\end{tabular}
\end{table}

\begin{figure}[h]
\centering
\includegraphics[width=0.62\textwidth]{figures/003nasionmap.jpg}
\caption{Eight Nasion polarity positions distributed symmetrically
  left and right of the facial midline.
  \textbf{Two positions of the same colour form one constitutional axis};
  polarity is always identical on both sides
  (e.g.\ Axis~1: both left and right positions are N for Pulmotonia,
  both S for Hepatonia).
  The midline itself is non-polar.
  Note that two adjacent positions on the right belong to
  \emph{different} axes --- do not conflate them.
  Blue: Axis~1 (1N/1S); Red: Axis~3 (3N/3S);
  Green: Axis~2 (2N/2S); Black: Axis~4 (4N/4S).}
\label{fig:nasionmap}
\end{figure}

Each pair sharing the same Nasion axis forms a constitutional
\emph{axis} with opposite polarities at the same two positions:
Axis~1 (1S$\leftrightarrow$1N), Axis~2 (2N$\leftrightarrow$2S),
Axis~3 (3N$\leftrightarrow$3S), Axis~4 (4N$\leftrightarrow$4S).

A single Nasion locus reading determines both axis (location) and
polarity (constitutional identity) simultaneously.  The 45-degree
rotational sweep of an east--west oriented probe magnet additionally
reveals constitutional-type-specific polarity reversal patterns,
providing supplementary identification information.\footnote{%
  The Nasion acupoint classification, constitutional axis structure,
  BIMOD detection protocol, and all related frameworks are the
  author's (DuckJin Yoon) independent hypotheses and observations.
  They are unrelated to the official Eight-Constitution Medicine of
  \url{www.ecmed.org}.  The theoretical motivation originates from
  the author's blog (``Eight Constitution and Yongshin Five-Element
  Hypothesis'') and is detailed in the BIMOD~001--005 paper series.}

% ─────────────────────────────────────────
\section{Supporting Evidence}

\subsection{Darkened Particle--S-Polarity Coupling (Novel Observation)}

\begin{figure}[h]
\centering
\includegraphics[width=0.72\textwidth]{figures/Magnetite_magnetosomes_SS5.png}
\caption{TEM image of magnetite magnetosomes in \textit{Gammaproteobacteria}
  strain SS-5 (panels A, B, C).
  Darkened/electron-dense particles correspond to S-polarity positions
  in Table~\ref{tab:tem} without exception (BIMOD analysis).
  \textit{Source:} Lefèvre et al.\ (2012), \textit{ISME J} 6:440--450,
  via Wikimedia Commons.
  Licence: CC~BY~3.0
  (\url{https://en.wikipedia.org/wiki/File:Magnetite_magnetosomes_in_Gammaproteobacteria.png}).}
\label{fig:ss5tem}
\end{figure}

\begin{table}[h]
\centering
\caption{BIMOD polarity readings: SS-5 magnetosome TEM images}
\label{tab:tem}
\begin{tabular}{p{2.2cm} p{5.8cm} p{4.5cm}}
\toprule
Panel & Polarity sequence (darkened in parentheses) & Observation \\
\midrule
(A) Full chain &
  SN(S)NSNSN(S)NSN(S)N &
  Darkened = S position \\[4pt]
(B) Partial, 8 particles &
  SNSNSN(S)N &
  Darkened = S position \\[4pt]
(C) Vertical view &
  N--S \ (box: N) &
  N terminal confirmed \\
\bottomrule
\end{tabular}
\end{table}

In TEM imaging, elevated electron density reflects higher Fe$^{3+}$
concentration or more complete oxidation (magnetite Fe$_3$O$_4$
$\to$ maghemite $\gamma$-Fe$_2$O$_3$).  The consistent co-occurrence
of darkened appearance and S-polarity assignment motivates the
following hypothesis:

\begin{quote}
\textit{S-polarity SMD BMNP nodes are maintained in a preferentially
oxidised or electron-dense state relative to N-polarity nodes,
generating a visually detectable contrast in TEM that parallels
their magnetic distinction.}
\end{quote}

If confirmed across additional specimens, darkened TEM particles
provide a morphological marker for BMNP polarity mapping without
BIMOD measurement.

\subsection{Single-Domain Confirmation}
\cite{kornig2020} confirms that each magnetosome is an independent
SMD nanoparticle exhibiting uniaxial anisotropy along the chain axis.
Magnetization reversal occurs only by coherent internal rotation.
This is structurally consistent with the BIMOD observation of
uniform single-polarity response across the entire nanoparticle area.

\subsection{Johng et al.\ (2007): Meridian-Pathway Transport}
Fluorescent magnetic nanoparticles injected at acupoint LR9 were
subsequently detected at LR3 on the same meridian; non-acupoint
injections showed no comparable distal migration \cite{johng2007}.
Note: injected foreign particles likely disrupted the native polarity
signal network at the injection site.  The finding confirms the
existence of a physical conduit following classical meridian routes.

\subsection{Plant BMNP Chains and Cross-Species Polarity Correspondence}
\cite{gorobets2023} identified BMNP chains in phloem sieve-tube cell
walls of tobacco, potato, and pea, following vascular conducting
pathways.  This suggests the cross-polar SMD network architecture is
a conserved biological strategy for long-range signal integration,
predating the nervous system.  If plant specimens yield constitutional
BIMOD readings, a polarity-correspondence principle between plant and
animal constitution---governing dietary affinity and antagonism---may
be physically grounded.

\subsection{East--West Geomagnetic Component and 45-Degree Dependency}
\cite{wang2019} demonstrated that alpha-wave suppression in response
to geomagnetic field rotation occurs when the east--west horizontal
component is independently controlled.  \cite{chae2019} showed
orientation toward east under blue-light fasting conditions.  These
findings confirm the east--west component as a biologically active
independent signal, providing physical context for the BIMOD
observation of constitutional-type-specific polarity reversals at
45-degree probe rotation intervals.

% ─────────────────────────────────────────
\section{Oxidation--Polarity Coupling Hypothesis}

\subsection{Physical Basis}
Magnetite (Fe$_3$O$_4$) contains both Fe$^{2+}$ and Fe$^{3+}$.
Surface oxidation to maghemite ($\gamma$-Fe$_2$O$_3$) converts all
iron to Fe$^{3+}$, changing the magnetic moment.  Higher Fe$^{3+}$
concentration appears darker in TEM.

\textbf{Proposed:} S-polarity SMD nodes are systematically maintained
in a more oxidised state than N-polarity nodes.  Possible causes:
(1)~differential access to metabolic electron donors along the chain;
(2)~directional electron flow under geomagnetic influence creating a
redox gradient; (3)~constitutional differences in iron-metabolism
enzymes producing genotype-specific oxidation profiles.

\subsection{Falsifiable Predictions}
\begin{itemize}
  \item In TEM images of BMNP chains from any species, darkened
        particles consistently occupy even-numbered (S-polarity) positions.
  \item XPS of isolated BMNP chains shows higher Fe$^{3+}$/Fe$^{2+}$
        at S-polarity positions.
  \item M\"{o}ssbauer spectroscopy detects isomer-shift differences
        between alternating particles.
  \item BIMOD assigns darkened TEM particles to S-polarity regardless
        of species or tissue source.
\end{itemize}

% ─────────────────────────────────────────
\section{Constitutional Pharmacology}

The organ dominance sequence unique to each constitution implies
differential expression of metabolic enzymes.  A critical extension is
proposed:

\begin{quote}
\textit{In constitutional axis pairs (natural mates in $\sim$99\% of
freely chosen couples), the therapeutically effective dose for one
polarity type corresponds to the adverse-effect dose for the opposite
polarity type---an extreme reciprocal relationship.}
\end{quote}

Example: Cholecystonia~(2N) vs.\ Colonotonia~(2S).
This reciprocal pharmacological principle may underpin optimal
personalised prescribing for drugs, vaccines, and therapeutic agents.

Specific hypotheses:
\begin{itemize}
  \item CYP2C9 (ibuprofen metabolism): higher activity in
        Hepatonia~1S, Cholecystonia~2N, Vesicotonia~3S.
  \item CYP2E1 (acetaminophen metabolism): higher activity in
        Pulmotonia~1N, Colonotonia~2S, Gastrotonia~3N, Renotonia~4S.
  \item Effect--adverse-effect inversion predicted between axis-paired
        constitutional types for the above drugs.
\end{itemize}

% ─────────────────────────────────────────
\section{Open Questions and Research Priorities}

\begin{table}[h]
\centering
\caption{Research priorities}
\begin{tabular}{p{5cm} p{5cm} l}
\toprule
Question & Method & Priority \\
\midrule
Do darkened TEM particles consistently occupy even (S) positions? &
  BIMOD analysis of additional published magnetosome TEM images &
  Immediate \\
Is each BMNP an SMD monopolar particle or a dipole? &
  SQUID / atomic magnetic microscopy on single particles &
  High \\
Which signal mechanism dominates: cross-polar relay vs.\ redox cascade? &
  Real-time SQUID along acupoint pathways &
  Medium \\
Do plant specimens yield BIMOD constitutional readings? &
  BIMOD protocol on plant tissue photographs &
  High \\
Can XPS / M\"{o}ssbauer confirm Fe oxidation at S-polarity positions? &
  Materials science collaboration &
  Medium \\
Does 45-degree sweep produce constitutional-specific reversals in
blinded trials? &
  Controlled angular sweep with typed subjects &
  High \\
Do axis-paired mates show 99\% frequency in free-choice couples? &
  Population-scale BIMOD constitutional typing of couples &
  High \\
\bottomrule
\end{tabular}
\end{table}

% ─────────────────────────────────────────
\section{Conclusion}
The BMNP Meridian Hypothesis positions the cross-polar SMD nanoparticle
chain---with monopolar terminal nodes---as the physical substrate of the
classical meridian system.  The novel observation that darkened TEM
particles consistently occupy S-polarity positions provides a potential
morphological marker for BMNP polarity identity.  The framework
integrates constitutional type, pharmacological response, mate
selection, and environmental adaptation under a single physical
principle: the human individual, constitutionally understood, is a
magneto-optical interferometer tuned by evolution to Earth's geomagnetic
field.  These insights constitute a scientific foundation for the era of
personalised medicine and constitutional dynamics in humanity.

% ─────────────────────────────────────────
\section*{Acknowledgements}
The author thanks twelve years of empirical observation conducted under
the principles of gratuitous service and resonant openness.  Collaborative
theoretical development was conducted with Claude~AI (Anthropic),
DeepSeek~AI, Gemini~AI (Google), and other AI systems.  The
Eight-Constitutional Research Platform community is thanked for sustained
observation and discussion.

% ─────────────────────────────────────────
\begin{thebibliography}{99}

\bibitem{yoon2026series}
Yoon,~D. (2026).
\textit{BIMOD 001--005: Foundational papers on the BIMOD framework}.
Eight-Constitutional Research Platform.
OSF: \url{https://doi.org/10.17605/OSF.IO/VJF63}

\bibitem{yoon2026_500}
Yoon,~D. (2026).
\textit{BIMOD-500: An Integrated Theory of Constitutional Classification
and Applications Based on Biogenic Magnetic Nanoparticles}.
Eight-Constitutional Research Platform.

\bibitem{johng2007}
Johng,~H.\,M.\ et al.\ (2007).
Use of magnetic nanoparticles to visualize threadlike structures inside
lymphatic vessels of rats.
\textit{Evid.\ Based Complement.\ Alternat.\ Med.}, 4(1), 77--82.

\bibitem{gorobets2019}
Gorobets,~O.\ et al.\ (2019).
Biogenic magnetic nanoparticles in Atlantic salmon olfactory tissue.
\textit{Scientific Reports}.

\bibitem{gorobets2023}
Gorobets,~O.\ et al.\ (2023).
Biogenic magnetic nanoparticles in plant phloem sieve-tube cell walls.
\textit{Plant and Soil}.

\bibitem{kornig2020}
K\"{o}rnig,~A.\ et al.\ (2020).
Stability and magnetic properties of magnetosome chains in
freeze-dried magnetotactic bacteria.
\textit{Nanoscale Advances}.

\bibitem{wang2019}
Wang,~C.\,X.\ et al.\ (2019).
Transduction of the geomagnetic field as evidenced from alpha-band
activity in the human brain.
\textit{eNeuro}, 6(2).

\bibitem{chae2019}
Chae,~K.\,S.\ et al.\ (2019).
Human magnetic sensing and orientation toward magnetic north and east.
\textit{PLoS ONE}.

\bibitem{wikipedia_magnetoreception}
Wikipedia.
\textit{Magnetoreception}.
\url{https://en.wikipedia.org/wiki/Magnetoreception}
(accessed February 2026).

\bibitem{ecmed}
Eight-Constitution Medical Research Institute.
\textit{Constitutional classification and official nomenclature}.
\url{www.ecmed.org}

\end{thebibliography}

\end{document}
